\section{Low-degree denominator computation table}
\label{sec:denominator-low-degree-computation-table}

\providecommand{\Z}{\mathbb Z}
\providecommand{\R}{\mathbb R}
\providecommand{\La}{\Lambda}
\providecommand{\autcor}{\mathfrak g_{\Delta_5}}
\providecommand{\sdim}{\operatorname{sdim}}
\providecommand{\smult}{\operatorname{smult}}
\providecommand{\ad}{\operatorname{ad}}
\providecommand{\Prim}{\operatorname{Prim}}

\subsection{Claim attacked}

The attacked claim is that the low-degree table of the Borcherds
denominator algebra already gives the microscopic Hall/CoHA bracket
constants.  The computation below gives a sharper statement.  The
active support and the Borcherds presentation determine the listed
denominator roots, signed supermultiplicities, and the normalized
Chevalley--Serre constants.  They do not determine compact Hall
extension constants.

\subsection{First-principles conversion}

Write a positive denominator degree as
$$
\beta=c_1\delta_1+c_2\delta_2+c_3\delta_3,
\qquad c_i\in\Z_{\geq0}.
$$
The manuscript degree map is
$$
\alpha(n,l,m)=n\delta_1+m\delta_2+(n+m-l)\delta_3.
$$
Hence the charge attached to \(\beta\) is
$$
\gamma(\beta)=(n,l,m)=(c_1,\ c_1+c_2-c_3,\ c_2).
$$
With
$$
A=((\delta_i,\delta_j))_{i,j}=
\begin{pmatrix}
2&-2&-2\\
-2&2&-2\\
-2&-2&2
\end{pmatrix},
$$
one computes
$$
(\beta,\beta)
=2(c_1^2+c_2^2+c_3^2)
-4(c_1c_2+c_1c_3+c_2c_3).
$$
For \(D(\gamma)=4nm-l^2\),
$$
(\beta,\beta)=-2D(\gamma).
$$
The denominator identity gives the signed root supermultiplicity
$$
\smult(\beta)=f(nm,l)
$$
for every active row \(f(nm,l)\neq0\).

\subsection{Jacobi coefficients used}

The coefficient computation of \(\phi_{0,1}\) through \(q^2\) gives
$$
f(0,-1)=1,\qquad f(0,0)=10,\qquad f(0,1)=1,
$$
$$
f(1,-2)=10,\quad f(1,-1)=-64,\quad f(1,0)=108,
\quad f(1,1)=-64,\quad f(1,2)=10,
$$
$$
f(2,-3)=1,\quad f(2,-2)=108,\quad f(2,-1)=-513,
\quad f(2,0)=808,\quad f(2,1)=-513,\quad
f(2,2)=108,\quad f(2,3)=1.
$$

\subsection{Active denominator rows through height three}

\begin{center}
\begingroup
\small
\begin{tabular}{c|c|c|c|c|c}
\(\operatorname{ht}\beta\) &
\(\beta\) &
\(\gamma(\beta)=(n,l,m)\) &
\((\beta,\beta)\) &
\(D=4nm-l^2\) &
\(f(nm,l)=\smult(\beta)\)\\
\hline
1 & \(\delta_1\) & \((1,1,0)\) & \(2\) & \(-1\) & \(f(0,1)=1\)\\
1 & \(\delta_2\) & \((0,1,1)\) & \(2\) & \(-1\) & \(f(0,1)=1\)\\
1 & \(\delta_3\) & \((0,-1,0)\) & \(2\) & \(-1\) & \(f(0,-1)=1\)\\
\hline
2 & \(\delta_1+\delta_2\) & \((1,2,1)\) & \(0\) & \(0\) & \(f(1,2)=10\)\\
2 & \(\delta_1+\delta_3\) & \((1,0,0)\) & \(0\) & \(0\) & \(f(0,0)=10\)\\
2 & \(\delta_2+\delta_3\) & \((0,0,1)\) & \(0\) & \(0\) & \(f(0,0)=10\)\\
\hline
3 & \(2\delta_1+\delta_2\) & \((2,3,1)\) & \(2\) & \(-1\) & \(f(2,3)=1\)\\
3 & \(\delta_1+2\delta_2\) & \((1,3,2)\) & \(2\) & \(-1\) & \(f(2,3)=1\)\\
3 & \(2\delta_1+\delta_3\) & \((2,1,0)\) & \(2\) & \(-1\) & \(f(0,1)=1\)\\
3 & \(\delta_1+2\delta_3\) & \((1,-1,0)\) & \(2\) & \(-1\) & \(f(0,-1)=1\)\\
3 & \(2\delta_2+\delta_3\) & \((0,1,2)\) & \(2\) & \(-1\) & \(f(0,1)=1\)\\
3 & \(\delta_2+2\delta_3\) & \((0,-1,1)\) & \(2\) & \(-1\) & \(f(0,-1)=1\)\\
3 & \(\delta_1+\delta_2+\delta_3\) & \((1,1,1)\) & \(-6\) & \(3\) & \(f(1,1)=-64\)
\end{tabular}
\endgroup
\end{center}

The inactive pure multiples in the same range are a useful sign check:
$$
\begin{array}{c|c|c|c|c}
\beta & \gamma(\beta) & (\beta,\beta)&D&f(nm,l)\\
\hline
2\delta_1&(2,2,0)&8&-4&f(0,2)=0\\
2\delta_2&(0,2,2)&8&-4&f(0,2)=0\\
2\delta_3&(0,-2,0)&8&-4&f(0,-2)=0\\
3\delta_1&(3,3,0)&18&-9&f(0,3)=0\\
3\delta_2&(0,3,3)&18&-9&f(0,3)=0\\
3\delta_3&(0,-3,0)&18&-9&f(0,-3)=0 .
\end{array}
$$

\subsection{Denominator bracket status}

Choose even Chevalley generators for the real simple roots:
$$
e_i\in\mathfrak g_{\delta_i},\qquad
f_i\in\mathfrak g_{-\delta_i},\qquad
h_i=[e_i,f_i].
$$
The normalized denominator constants are
$$
[h_i,h_j]=0,\qquad
[h_i,e_j]=A_{ij}e_j,\qquad
[h_i,f_j]=-A_{ij}f_j,\qquad
[e_i,f_j]=\delta_{ij}h_i.
$$
For \(i\neq j\), put
$$
E_{ij}:=[e_i,e_j]\in\mathfrak g_{\delta_i+\delta_j},
\qquad E_{ji}=-E_{ij}.
$$
If \(\{i,j,k\}=\{1,2,3\}\), then Jacobi gives
$$
[h_i,E_{ij}]=[h_j,E_{ij}]=0,\qquad
[h_k,E_{ij}]=-4E_{ij},
$$
$$
[f_i,E_{ij}]=2e_j,\qquad
[f_j,E_{ij}]=-2e_i,\qquad
[f_k,E_{ij}]=0.
$$
For ordered \(i\neq j\), put
$$
E_{iij}:=[e_i,E_{ij}]=[e_i,[e_i,e_j]]
\in\mathfrak g_{2\delta_i+\delta_j}.
$$
The real Serre relation is
$$
[e_i,E_{iij}]=(\ad e_i)^3e_j=0,
$$
and the lowering constants are
$$
[f_i,E_{iij}]=2E_{ij},\qquad
[f_j,E_{iij}]=0,\qquad
[f_k,E_{iij}]=0.
$$

For \(i<j\), the primitive isotropic degree
$$
a_{ij}:=\delta_i+\delta_j
$$
has signed denominator multiplicity \(10\).  One even direction is the
real-real commutator \(E_{ij}\).  The correction character gives nine
additional even isotropic simple generators on the same primitive ray:
$$
u_{ij,1},\ldots,u_{ij,9}\in\mathfrak g_{a_{ij}},
\qquad \tau(ta_{ij})=9\quad(t\geq1),
$$
so
$$
\smult(a_{ij})=1+9=10.
$$
The orthogonality relations are
$$
[e_i,u_{ij,r}]=[e_j,u_{ij,r}]=0,\qquad
[f_i,u_{ij,r}]=[f_j,u_{ij,r}]=0,
$$
and the non-orthogonal real root only forces
$$
(\ad e_k)^5u_{ij,r}=0.
$$
Thus
$$
[e_k,u_{ij,r}]
\in\mathfrak g_{\delta_1+\delta_2+\delta_3}
$$
is not a numerical structure constant until a basis of that root space
has been constructed.

In degree \(\delta_{123}:=\delta_1+\delta_2+\delta_3\), set
$$
T_1=[e_1,E_{23}],\qquad
T_2=[e_2,E_{31}],\qquad
T_3=[e_3,E_{12}].
$$
Jacobi gives
$$
T_1+T_2+T_3=0.
$$
The denominator product gives only the signed total
$$
\smult(\delta_{123})=f(1,1)=-64.
$$
It does not choose a basis of
\(\mathfrak g_{\delta_{123}}\), split even and odd dimensions, or give
the constants expressing \([e_k,u_{ij,r}]\) in such a basis.

\subsection{Replacement tables for \texttt{proj.tex}}

The following two tables are the exact low-degree replacement material.
They separate denominator-algebra facts from the unproved Hall
realization.

\begin{center}
\begingroup
\small
\begin{tabular}{c|c|c|c|c|c}
\(\operatorname{ht}\) & active charge \(\gamma\) & root \(\beta\) &
norm & discriminant & expected \(\smult\)\\
\hline
1 & \((1,1,0)\) & \(\delta_1\) & \(2\) & \(-1\) & \(1\)\\
1 & \((0,1,1)\) & \(\delta_2\) & \(2\) & \(-1\) & \(1\)\\
1 & \((0,-1,0)\) & \(\delta_3\) & \(2\) & \(-1\) & \(1\)\\
2 & \((1,2,1)\) & \(\delta_1+\delta_2\) & \(0\) & \(0\) & \(10\)\\
2 & \((1,0,0)\) & \(\delta_1+\delta_3\) & \(0\) & \(0\) & \(10\)\\
2 & \((0,0,1)\) & \(\delta_2+\delta_3\) & \(0\) & \(0\) & \(10\)\\
3 & \((2,3,1)\) & \(2\delta_1+\delta_2\) & \(2\) & \(-1\) & \(1\)\\
3 & \((1,3,2)\) & \(\delta_1+2\delta_2\) & \(2\) & \(-1\) & \(1\)\\
3 & \((2,1,0)\) & \(2\delta_1+\delta_3\) & \(2\) & \(-1\) & \(1\)\\
3 & \((1,-1,0)\) & \(\delta_1+2\delta_3\) & \(2\) & \(-1\) & \(1\)\\
3 & \((0,1,2)\) & \(2\delta_2+\delta_3\) & \(2\) & \(-1\) & \(1\)\\
3 & \((0,-1,1)\) & \(\delta_2+2\delta_3\) & \(2\) & \(-1\) & \(1\)\\
3 & \((1,1,1)\) & \(\delta_1+\delta_2+\delta_3\) & \(-6\) & \(3\) & \(-64\)
\end{tabular}
\endgroup
\end{center}

\begin{center}
\begingroup
\small
\begin{tabular}{p{0.19\textwidth}|p{0.38\textwidth}|p{0.33\textwidth}}
\textbf{Degree} & \textbf{Denominator bracket status} & \textbf{Hall status}\\
\hline
\(\delta_i\) &
Real simple generator \(e_i\); Chevalley constants fixed by \(A\). &
Requires a compact primitive class in charge \(\gamma_i\); not supplied by the determinant.\\
\hline
\(\delta_i+\delta_j\) &
\(E_{ij}=[e_i,e_j]\) plus nine isotropic simple directions
\(u_{ij,r}\); \(10=1+9\). &
Requires the oriented compact extension correspondence for the two
simple charges; determinant gives only signed size \(10\).\\
\hline
\(2\delta_i+\delta_j\) &
\(E_{iij}=[e_i,[e_i,e_j]]\), with
\((\ad e_i)^3e_j=0\) and
\([f_i,E_{iij}]=2E_{ij}\). &
Requires Hall Serre and lowering constants; not derived from the Euler
pairing or the Fock determinant.\\
\hline
\(\delta_1+\delta_2+\delta_3\) &
Real bracketings \(T_1,T_2,T_3\) satisfy \(T_1+T_2+T_3=0\);
\(\smult=-64\).  Constants involving \([e_k,u_{ij,r}]\) need a basis. &
First fatal Hall gap: no compact correspondence, orientation cocycle,
protected integration, basis, or parity split has been constructed.
\end{tabular}
\endgroup
\end{center}

\subsection{Why the Hall constants are not derivable here}

A compact Hall derivation would have to start from extension
correspondences
$$
\mathfrak M_\gamma^\sigma\times\mathfrak M_\eta^\sigma
\xleftarrow{\ p_{\gamma,\eta}\ }
\mathfrak E_{\gamma,\eta}^\sigma
\xrightarrow{\ q_{\gamma,\eta}\ }
\mathfrak M_{\gamma+\eta}^\sigma
$$
with vanishing cycles and an orientation line, producing
$$
m_{\gamma,\eta}
=q_{\gamma,\eta!}p_{\gamma,\eta}^{*}
$$
and then a protected primitive supercommutator
$$
[u_{\gamma,a},u_{\eta,b}]_{\mathrm{Hall}}
=
\sum_c C^c_{(\gamma,a),(\eta,b)}u_{\gamma+\eta,c}.
$$
The denominator product supplies only
$$
\sum_a(-1)^{|a|}
=\sdim\Prim_{\mathrm{prot},\gamma}
=f(nm,l).
$$
It supplies no stacks \(\mathfrak M_\gamma^\sigma\), no compact
extension stack \(\mathfrak E_{\gamma,\eta}^\sigma\), no orientation
cocycle, no pull-push class, and no basis of protected primitives.
Therefore it supplies no integers
\(C^c_{(\gamma,a),(\eta,b)}\).

The determinant cannot recover the missing constants.  Any graded super
vector space with the same signed dimensions gives the same product, and
the zero bracket on that graded space gives the same determinant as the
Borcherds bracket.  Since the Borcherds denominator algebra has
\([e_i,e_j]=E_{ij}\neq0\) for \(i\neq j\), the determinant is strictly
weaker than the bracket.

\subsection{Verification record}

The lattice calculation checks the Gram matrix
$$
((\delta_i,\delta_j))=
\begin{pmatrix}
2&-2&-2\\
-2&2&-2\\
-2&-2&2
\end{pmatrix},
$$
the Weyl vector
$$
\rho=\frac12(\delta_1+\delta_2+\delta_3),
$$
and the anti-isometry
$$
(\alpha(n,l,m),\alpha(n,l,m))=-2(4nm-l^2).
$$
The Jacobi-form calculation checks the coefficients of \(\phi_{0,1}\)
through \(q^2\) and the eta-product coefficient list for
\(\prod_{k\geq1}(1-q^k)^9\).  A one-off height enumeration produced all
active rows in the displayed table and the inactive pure multiples.

\subsection{Remaining obligations}

The denominator table is closed through height three.  The microscopic
Hall table is not closed.  The missing inputs are:
$$
(\mathfrak M_\gamma^\sigma,\mathfrak E_{\gamma,\eta}^\sigma,
o,\varepsilon_{\gamma,\eta},\mathrm{Int}_{\mathrm{prot}})
$$
for the compact \(K3\times E\) problem, with proof that the protected
primitive supercommutator satisfies the Chevalley, Serre,
orthogonality, and imaginary-root relations of \(\autcor\).  In degree
\(\delta_1+\delta_2+\delta_3\) one must additionally construct a basis
and parity split of the protected primitive space and compute the
constants expressing \([e_k,u_{ij,r}]\) in that basis.
